% Retained detailed derivation for the complete integral-value certificate.
% This is a supplementary source fragment, not a replacement manuscript.
\subsection{Complete arithmetic and continuous-price certificates}
\label{sec:integral-full-certificate}

I state the finite-law, scalar-arithmetic and monetary error premises separately,
because a unitary circuit which evaluates a finite table correctly does not,
by itself, price the original continuous financial claim accurately.  The
arguments below certify the particular integral and direct constructions.  They
do not establish that either construction is the cheapest possible quantum or
classical implementation.

Let $K>0$, $H>0$ and $\lambda>0$.  In the retained implementation $K=100$.
Write $\phi$ and $\Phi$ for the standard normal density and distribution
function.  The conditional capped-payoff kernel is
\begin{equation}
 P(m)=\int_K^{K+H}\Phi\!\left(\frac{m-\log(k/K)}{\lambda}\right)\,dk
     =H\int_0^1\Phi\!\left(\frac{m-\log(1+Hs/K)}{\lambda}\right)\,ds.
 \label{eq:iv-P-strike}
\end{equation}
This equality follows by expressing the capped payoff as the integral of cash
digitals and then applying Tonelli's theorem to its nonnegative integrand.
In particular $0\leq P\leq H$.  Differentiation under the integral is justified
at every finite order by the boundedness of a polynomial times $\phi$.
For $j\geq1$,
\begin{equation}
 P^{(j)}(m)=\int_K^{K+H}\lambda^{-j}(-1)^{j-1}
 \operatorname{He}_{j-1}\!\left(\frac{m-\log(k/K)}{\lambda}\right)
 \phi\!\left(\frac{m-\log(k/K)}{\lambda}\right)\,dk.
 \label{eq:iv-P-derivative}
\end{equation}
Here $\operatorname{He}_j$ denotes the probabilists' Hermite polynomial.
If $M_j\geq\sup_x|\phi^{(j)}(x)|$ is a retained upper certificate, then
$\|P^{(j)}\|_\infty\leq HM_{j-1}\lambda^{-j}$.  Positivity gives $P'\geq0$.
The pathwise derivative of the capped payoff with respect to its log-scale is
$Ke^{m+\lambda Z}\boldsymbol 1_{K<Ke^{m+\lambda Z}<K+H}$ except at its
two thresholds, which have probability zero.  This derivative is bounded by
$K+H$.  The payoff as a function of its log-scale is itself $(K+H)$-Lipschitz,
so dominated convergence for difference quotients gives $P'\leq K+H$.
I therefore use
\begin{equation}
 P'_* = \min\{K+H,HM_0/\lambda\}.
 \label{eq:iv-Pprime-bound}
\end{equation}

Let $w_i>0$, $\sum_iw_i=1$, $d_i\in\mathbb R$ and
$s_i\in\mathbb R^p$ satisfy
\begin{equation}
 \sum_iw_id_i=0,\qquad \sum_iw_is_i=0.
 \label{eq:iv-centre}
\end{equation}
Define $T(y)=\sum_iw_i\exp(d_i+s_i\cdot y)$,
$U(y)=T(y)-1$ and $m(y)=m_0+a\cdot y$.  Jensen's inequality and
\eqref{eq:iv-centre} give $T(y)\geq1$ for every $y$, so $U(y)\geq0$.
These are the same geometric-basket ratio and common-factor conditioning
coordinates as in the classical risk-neutral model.  No change of financial
dynamics is made here.

Put $g_m(t)=P(m+\log t)$ for $t>0$ and
\begin{equation}
 A_2(x)=-\lambda^{-1}-x\lambda^{-2},\qquad
 D_2\geq\sup_{x\in\mathbb R}|A_2(x)\phi(x)|.
 \label{eq:iv-A2}
\end{equation}
The retained rational $D_2$ is an outward upper endpoint of the previously
certified joint Gaussian derivative norm.  Direct differentiation of
\eqref{eq:iv-P-strike} gives
\begin{equation}
 g_m''(t)=\frac{H}{t^2}\int_0^1
 A_2\!\left(\frac{m+\log t-\log(1+Hs/K)}{\lambda}\right)
 \phi\!\left(\frac{m+\log t-\log(1+Hs/K)}{\lambda}\right)\,ds.
 \label{eq:iv-g2}
\end{equation}
Taylor's theorem in integral form, with the substitution $t=1+zU$, gives
\begin{align}
 R_2(y)&=P(m(y)+\log T(y))-P(m(y))-U(y)P'(m(y))\notag\\
 &=\frac{HU(y)^2}{2}\mathbb E_{Z,S}
 \left[\frac{A_2(\xi)\phi(\xi)}{(1+ZU(y))^2}\right],
 \quad \xi=\frac{m(y)+\log(1+ZU(y))-\log(1+HS/K)}{\lambda},
 \label{eq:iv-R2-identity}
\end{align}
where $S$ is uniform on $[0,1]$ and $Z$ has density $2(1-z)$ on $[0,1]$.
One exact construction of $Z$ is the minimum of two independent uniforms.
Indeed $\mathbb P(\min(V_1,V_2)>z)=(1-z)^2$.
The bounded first correction is
\begin{equation}
 R_1(y)=P(m(y)+\log T(y))-P(m(y)),\qquad
 0\leq R_1(y)\leq \min\{H,HD_1U(y)\},
 \label{eq:iv-R1}
\end{equation}
where $D_1\geq\sup_x\lambda^{-1}\phi(x)$.

\subsubsection{Finite preparation and admissible normalisation}

For a finite quantile construction let $L>0$, $N=2^{b_y}$,
$\alpha_1=\Phi(L)-\Phi(-L)$ and
\begin{equation}
 y_j=\Phi^{-1}\!\left(\Phi(-L)+\alpha_1\frac{j+1/2}{N}\right),
 \qquad 0\leq j<N.
 \label{eq:iv-quantile}
\end{equation}
An outward enclosure of each $y_j$ is computed, and the actual circuit
coordinate $\widehat y_j$ is the retained integer divided by $2^q$, minus the
specified offset $L$.  All comparisons and finite moments below use these
actual dyadic coordinates.  They are not silently replaced by the ideal
quantiles.  The table is checked to be nondecreasing and to lie in $[-L,L]$.
Independent $b_y$-qubit Hadamard registers prepare all $N^p$ coordinate labels
with equal amplitudes.  A reversible integer lookup computes their coordinates.
Duplicate rounded coordinates keep their distinct labels and their weights.
The retained Gaussian box factor is $\alpha=\alpha_1^p$.

A second preparation uses independent uniform coordinate midpoints
$y_j=-L+2L(j+1/2)/N$, with a dyadic $L$ and a large enough $q$ for exact
representation.  Its relative density weight is
$\omega(y)=\exp(-\|y\|^2/2)\in[0,1]$, and its external monetary factor is
\begin{equation}
 C_L=(2L/\sqrt{2\pi})^p.
 \label{eq:iv-box-factor}
\end{equation}
Thus $C_L\mathbb E_{\mathrm{uniform\ box}}[\omega(y)R(y)]$ is the
continuous Gaussian box integral, before midpoint discretisation.
Neither construction assumes a free Gaussian state-preparation oracle.

A finite normaliser must dominate $U$ at every realised coordinate, but it
need not dominate the discarded continuous tail.  Let $[\ell,h]$ cover the
outward quantile-node enclosures and the actual dyadic coordinates.  Since
$T$ is convex in each coordinate separately,
\begin{equation}
 U_*\geq\max_{v\in\{\ell,h\}^p}T(v)-1
 \label{eq:iv-vertex}
\end{equation}
is valid throughout the tensor box.  To prove this, fix all but one
coordinate and bound the value of its convex univariate function by the
larger endpoint value.  Apply this argument successively to all $p$
coordinates.  Each of the $2^p$ vertex sums is evaluated outwards, so the
maximum of their upper endpoints is a valid bound.  This is generally tighter
than adding the separate peak of every exponential term.  The full
continuous tail and grid transfer remain separate error premises.

If $U_*=0$, the correction vanishes and no correction estimator is required.
Otherwise set
\begin{equation}
 B_2=HD_2U_*^2/2,\qquad B_1=\min\{H,HD_1U_*\}.
 \label{eq:iv-normalisers}
\end{equation}
The same $B_2$ applies to the integral and direct second corrections, because
\eqref{eq:iv-R2-identity} implies
$|R_2(y)|\leq B_2U(y)^2/U_*^2$.
Writing
\begin{equation}
 a(y,z)=\left(\frac{U(y)}{U_*}\right)^2(1+zU(y))^{-2},\qquad
 h(x)=\frac{1+A_2(x)\phi(x)/D_2}{2},
 \label{eq:iv-a-h}
\end{equation}
gives $0\leq a,h\leq1$.  The integral circuit computes
\begin{equation}
 p_2(y,z,s)=\frac{1-a(y,z)}2+a(y,z)h(\xi)
          =\frac{1+X_2(y,z,s)}2,
 \qquad |X_2|\leq U(y)^2/U_*^2.
 \label{eq:iv-probability}
\end{equation}
The direct second-correction probability is
$\{1+R_2(y)/B_2\}/2$; the positive first-correction probability is
$R_1(y)/B_1$.  With a uniform-box preparation, $X_2$ is multiplied by
$\omega$, or the positive first-correction probability is multiplied by
$\omega$.  Final projection onto $[0,1]$ ensures a valid comparator threshold
and is nonexpansive relative to the ideal probability already in that interval.

\subsubsection{Certified scalar tables}

Let $f\in C^2([a,a+s])$, with $s>0$ a power of two.  Its endpoint grid is
$x_j=a+sj/M$, $0\leq j\leq M$, where $M=2^b$ and
$b\leq q+\log_2s$.  Suppose the retained unsigned integer $n_j$ satisfies
\begin{equation}
 |2^{-q}n_j-f(x_j)|\leq e_j.
 \label{eq:iv-node-enclosure}
\end{equation}
The implementation stores $n_j$, $|n_{j+1}-n_j|$ and the decreasing bit
$\boldsymbol 1_{n_{j+1}<n_j}$.  It multiplies the unsigned magnitude by
the within-cell integer remainder, drops the appropriate lower bits, then
adds or subtracts the resulting nonnegative adjustment according to that bit.
The final grid endpoint is stored explicitly, so an argument exactly at the
right endpoint has its own address.  No modular representation of a negative
slope is multiplied.

Let $L_1\geq\sup|f'|$, $L_2\geq\sup|f''|$ and $e_*\geq\max_je_j$.
For every represented argument in the declared interval,
\begin{equation}
 |\widehat f(x)-f(x)|\leq e_*+\frac{L_2s^2}{8M^2}+2^{-q}.
 \label{eq:iv-interpolation}
\end{equation}
To see this, let $I f$ be the exact linear interpolant on a cell
$[u,v]$.  The two-endpoint interpolation error is
$f(x)-I f(x)=f''(\zeta)(x-u)(x-v)/2$ for some $\zeta\in(u,v)$.
This follows by subtracting the linear interpolant, subtracting a multiple
of $(t-u)(t-v)$ chosen to give a third zero at $x$, and applying Rolle's
theorem twice.  Since $(x-u)(v-x)\leq(v-u)^2/4$, its magnitude is at most
$L_2(v-u)^2/8$.  Interpolating the two rounded endpoint values adds at most
$e_*$, by their nonnegative interpolation weights.  Integer truncation of
the adjustment adds less than one output unit $2^{-q}$, whether the table
is increasing or decreasing.  This proves \eqref{eq:iv-interpolation}.
For an argument with absolute error at most $\eta_x$, add $L_1\eta_x$,
provided both exact and implemented arguments are in the declared interval.
A preliminary projection of the implemented input onto that interval does
not enlarge its distance from an admissible exact argument.

Every endpoint value is enclosed using outward Arb arithmetic.  For portable
metadata, an exact rational lower endpoint is rounded down and an upper
endpoint up to the lattice $2^{-384}\mathbb Z$.  These enlarged intervals,
rather than the narrower internal balls, determine $e_j$ in
\eqref{eq:iv-node-enclosure}.  Nearest rounding to the output lattice is
chosen from their midpoint; the largest distance from the chosen dyadic value
to either enlarged endpoint is retained.  This also covers extreme Gaussian
tails without assuming that a printed floating-point zero is exact.

The whole-domain derivative bounds used in these tables are
\begin{align}
 e^x &: \quad L_1=L_2=e^{a+s},\notag\\
 \log x &: \quad L_1=a^{-1},\quad L_2=a^{-2},\quad a>0,\notag\\
 x^{-2} &: \quad L_1=2a^{-3},\quad L_2=6a^{-4},\quad a>0,\notag\\
 x^2 &: \quad L_1=2\max\{|a|,|a+s|\},\quad L_2=2,\notag\\
 h(x) &: \quad
 L_1=\frac{M_1/\lambda+M_2/\lambda^2}{2D_2},\quad
 L_2=\frac{M_2/\lambda+M_3/\lambda^2}{2D_2},\notag\\
 P(m) &: \quad L_1=P'_*,\quad L_2=HM_1/\lambda^2,\notag\\
 P'(m) &: \quad L_1=HM_1/\lambda^2,\quad L_2=HM_2/\lambda^3.
 \label{eq:iv-scalar-derivatives}
\end{align}
The first four follow by ordinary differentiation.  For $h$, use
$A_2\phi=-\phi/\lambda+\phi'/\lambda^2$ and then differentiate once or
twice.  The last two follow from \eqref{eq:iv-P-derivative} and
\eqref{eq:iv-Pprime-bound}.  The range checks use $h\in[0,1]$, $P\in[0,H]$
and $P'\in[0,P'_*]$.  They are analytical bounds on the full declared
interval, not sampled derivative maxima.

\subsubsection{Absolute arithmetic propagation and integer guards}

Each graph node carries an ideal real range, an independently propagated
implemented range, and an absolute error.  Suppose
$|\widehat x-x|\leq\eta_x$ and $|\widehat y-y|\leq\eta_y$.
Addition and subtraction have error at most $\eta_x+\eta_y$.  Projection
onto a fixed closed interval is nonexpansive.  Positive subtraction
$\max(0,x-y)$ therefore has the same error bound.  A rounded projection
endpoint introduces at most its own rounding error.

For a product followed by one fixed-point truncation, either identity
\begin{equation}
 \widehat x\widehat y-xy=x(\widehat y-y)+\widehat y(\widehat x-x)
 =y(\widehat x-x)+\widehat x(\widehat y-y)
 \label{eq:iv-product-identity}
\end{equation}
gives a valid upper bound.  If $X,Y,\widehat X,\widehat Y$ bound the absolute
values of the corresponding real quantities, the retained error is
\begin{equation}
 \eta_{xy}\leq
 \min\{X\eta_y+\widehat Y\eta_x,
        Y\eta_x+\widehat X\eta_y\}+2^{-q}.
 \label{eq:iv-product-error}
\end{equation}
For a nonnegative scale $c$ implemented as $\widehat c$,
\begin{equation}
 |\widehat c\widehat x-cx|\leq
 c\eta_x+\widehat X|\widehat c-c|,
 \label{eq:iv-scale-error}
\end{equation}
with one additional output unit when a right shift truncates the result.
Exact dyadic scales use their exact integer and shift; other scales use the
retained nearest $q$-bit coefficient.  Scalar-function errors are propagated
with \eqref{eq:iv-interpolation} and $L_1\eta_x$.  All these are absolute
bounds.  They do not divide by a small difference of capped-price terms.

Implemented interval endpoints are computed for the actual rounded
constants, table nodes, additions, differences, products and shifts.  They
may then be intersected with the independently justified ideal range
expanded by the absolute error.  This intersection retains correlations
such as a tightly bounded affine combination, without suppressing its
rounding error.  A contradictory intersection rejects the configuration.

Every unsigned operation is checked to have a nonnegative output below the
representable upper limit.  Positive subtraction uses an unsigned comparison
to copy the difference only when its operands are ordered; its temporary
modular difference is never fed into a financial multiplication.  The direct
second-correction construction forms separately nonnegative positive and
negative scaled terms.  Only their final difference uses two's-complement
arithmetic, followed immediately by an explicit signed projection onto
$[0,1]$.  Its full real interval must lie strictly within the signed
representable range, including all propagated errors.  A negative signed
intermediate is not subsequently treated as an unsigned multiplicand.
Endpoint table addressing, interpolation magnitudes, all scale constants and
final comparator thresholds are checked against their actual integer widths.
The final threshold lies in $\{0,\ldots,2^q\}$ and uses $q+1$ output bits.
A $q$-bit uniform comparator, with a zero-extended high bit, consequently has
success probability exactly equal to this threshold divided by $2^q$,
including both probability endpoints zero and one.

If the resulting probability error certificate is $\eta_p$, the monetary
value error before the external factor is at most $2B_2\eta_p$ for either
second correction and $B_1\eta_p$ for the positive first correction.
No continuous-tail, grid or estimator-confidence allowance is included in
this arithmetic certificate.

\subsubsection{Exact finite moments and two rare signed flags}

The moment certificate uses the realised financial labels alone.  It does
not require a financial price reference.  If $Y$ has an independent product
of the retained coordinate laws, define its one-dimensional moment generating
sum by
\begin{equation}
 \mathcal M(v)=\frac1N\sum_{j=0}^{N-1}\exp(v\widehat y_j).
 \label{eq:iv-finite-mgf}
\end{equation}
For a uniform-box law replace each summand by
$\exp(vy_j-y_j^2/2)$, thus including the relative density weight.
Expanding a finite power of $T$ gives
\begin{equation}
 \mathbb E[\omega(Y)T(Y)^r]
 =\sum_{|\gamma|=r}\frac{r!}{\prod_i\gamma_i!}
    \left(\prod_iw_i^{\gamma_i}\right)
    \exp\!\left(\sum_i\gamma_id_i\right)
    \prod_{k=1}^p\mathcal M\!\left(\sum_i\gamma_is_{ik}\right),
 \label{eq:iv-finite-raw-moment}
\end{equation}
where $\omega=1$ for quantiles and is the relative Gaussian weight for the
uniform box.  This follows from the multinomial theorem and independence;
all sums are finite.  The centred moments follow from
\begin{equation}
 \mathbb E[\omega(Y)U(Y)^r]
 =\sum_{j=0}^r(-1)^{r-j}\binom rj
   \mathbb E[\omega(Y)T(Y)^j].
 \label{eq:iv-finite-centred-moment}
\end{equation}
In particular the zeroth weighted moment is not replaced by one for a
uniform-box construction.  Every exponential sum and subsequent cancellation
is evaluated outwards.  This reduces a tensor-grid moment calculation to a
finite collection of one-dimensional sums, not $N^p$ price evaluations.
The calculation and its classical preprocessing cost remain recorded.

For a represented probability $\widehat p\in[0,1]$, define
\begin{equation}
 \widehat p_+=\max(0,2\widehat p-1),\qquad
 \widehat p_-=\max(0,1-2\widehat p).
 \label{eq:iv-rare-flags}
\end{equation}
Both are exactly represented on the same dyadic lattice, and
$\widehat p_+-\widehat p_-=2\widehat p-1$ exactly.  For the ideal integral
probability in \eqref{eq:iv-probability}, each ideal flag is at most
$U^2/U_*^2$, since $|2h-1|\leq1$ and $(1+zU)^{-2}\leq1$.
For the ideal direct second correction the same result follows from its
pointwise Taylor bound.  With a uniform-box preparation both are at most
$\omega U^2/U_*^2$.  Since $x\mapsto\max(0,x)$ is one-Lipschitz and
$|\widehat p-p|\leq\eta_p$, each implemented success probability has the
rigorous prior
\begin{equation}
 \mathbb E\widehat p_\pm\leq
 \min\!\left\{1,\frac{\mathbb E[\omega(Y)U(Y)^2]}{U_*^2}+2\eta_p\right\}.
 \label{eq:iv-rare-prior}
\end{equation}
The expectation here is under the prepared conditional quantile law or the
uniform label law.  Its external monetary factor $\alpha$ or $C_L$ does not
belong in this probability prior.  The combined arithmetic error is still
$2B_2\eta_p$, because the implemented flag difference is exact; it is not
necessary to charge two independent pointwise arithmetic estimates of that
difference.  Separate statistical estimates of the two expectations do,
however, require a joint confidence allocation.  A success-prior reduction
is not a physical-cost reduction until both complete flag preparations and
their chosen estimators have been charged.

The positive first-correction comparator also receives a finite-law mean
prior rather than the uninformative upper bound one.  The capped payoff is
one-Lipschitz in its underlying price and is nondecreasing.  Conditional on
$y$, the arithmetic underlying is $T(y)$ times the geometric underlying;
their expected positive difference is
$K\exp(m(y)+\lambda^2/2)U(y)$ after integrating the independent common
Gaussian factor.  Consequently
\begin{equation}
 0\leq\mathbb E[\omega R_1]\leq
 K e^{m_0+\lambda^2/2}\left[
 \sum_iw_ie^{d_i}\prod_{j=1}^p\mathcal M(a_j+s_{ij})
 -\prod_{j=1}^p\mathcal M(a_j)\right].
 \label{eq:iv-R1-correlated-gap}
\end{equation}
The finite sums $\mathcal M$ are those of
\eqref{eq:iv-finite-mgf}, including the relative density when appropriate.
This expansion follows by multiplying $T-1$ by $e^{a\cdot y}$ and applying
independence.  It retains the correlation between the level and the ratio;
replacing them by separate expectations would generally be invalid.
The pointwise bound in \eqref{eq:iv-R1} also yields
$\mathbb E[\omega R_1]\leq HD_1\mathbb E[\omega U]$, and boundedness gives
$\mathbb E[\omega R_1]\leq H\mathbb E\omega$.  I use the smallest of
these three outward upper bounds.  Dividing by $B_1$ and adding its own
probability arithmetic error $\eta_{p,1}$ gives a valid prior for the actual
first-correction success probability.  No price reference or external
monetary factor enters this finite prepared-law probability bound.

\subsubsection{Auxiliary and financial discretisation errors}

Couple each ideal uniform auxiliary variable to the midpoint of its dyadic
cell.  Their distance is at most $2^{-b-1}$.  The minimum map is
one-Lipschitz in the maximum norm, so the same bound applies to the minimum
of two coupled variables.  Let $D_3$ be the certified third joint norm and
$E_2\geq\sup_x|(A_2\phi)'(x)|$.  Differentiation of the Taylor integrand
with respect to $z$, before normalising it, gives magnitude at most
$HD_3U^3/2$.  Differentiation with respect to $s$ gives magnitude at most
$H^2E_2U^2/(2K\lambda)$, because
$|\partial_s\log(1+Hs/K)|\leq H/K$.
Under the actual finite financial law, the monetary auxiliary bias is therefore
\begin{equation}
 \beta_Z\leq\frac{HD_3\mathbb E[\omega U^3]}{2^{b_Z+2}},\qquad
 \beta_S\leq\frac{H^2E_2\mathbb E[\omega U^2]}
                         {K\lambda\,2^{b_S+2}}.
 \label{eq:iv-aux-bias}
\end{equation}
These are before the external factor.  Continuous Gaussian moments cannot
be substituted for finite-law moments without their own transfer bound.
A valid $E_2$ is $M_1/\lambda+M_2/\lambda^2$; the general signed-Stirling
triangle bound also applies.  This conservative estimate does not claim
that either auxiliary width is optimal.

For the one-dimensional retained prototypes, let $[V_-,V_+]$ be an
independently validated enclosure of the original continuous undiscounted
price $V$, and let $[c_-,c_+]$ enclose the continuous Taylor coefficient
sum $c_0$ or $c_0+c_1$.  Let $[J_-,J_+]$ enclose the exact finite,
mass-weighted correction, using the actual dyadic financial labels and exact
conditional kernels.  Then
\begin{equation}
 |J-(V-c)|\leq
 \max\{|J_--(V_+-c_-)|,\ |J_+-(V_--c_+)|\}.
 \label{eq:iv-reference-transfer}
\end{equation}
Indeed $J-(V-c)$ lies between the two displayed endpoint differences.  This
bound includes the discarded continuous tail, ideal financial grid and
financial-node rounding together.  It does not include auxiliary sampling,
value-oracle arithmetic or statistical estimation.  The source interval,
reference hash, actual source version and finite computation are retained.
The continuous reference is paid validation or preprocessing: it already
resolves the price, and cannot be treated as free information in a practical
quantum-advantage comparison.

For a higher-dimensional analytical alternative, let $R_r$ denote the
order-$r$ correction, $r=1$ or $2$, and let $Q(z)$ be the Taylor integrand
$g_m^{(r)}(1+zU)/H$.  On a full box, suppose $U\leq\overline U$,
$|\partial_jU|\leq U_j$ and $|\partial_{jj}U|\leq U_{jj}$.
Let $D_r,D_{r+1},D_{r+2}$ bound the relevant $t$ derivatives, and
$E_r,E_{r+1},F_r$ bound the first $x$ derivative, mixed derivative and
second $x$ derivative of their Gaussian kernels.  The integral remainder
and the chain rule imply
\begin{align}
 |\partial_jR_r|&\leq\frac H{r!}
 \left[r\overline U^{r-1}U_jD_r+
 \overline U^r\left(|a_j|E_r/\lambda+U_jD_{r+1}\right)\right],
 \label{eq:iv-R-first}\\
 |\partial_{jj}R_r|&\leq\frac H{r!}\left[
 r(r-1)\overline U^{r-2}U_j^2D_r+
 r\overline U^{r-1}U_{jj}D_r\right.\notag\\
 &\quad+2r\overline U^{r-1}U_j
           (|a_j|E_r/\lambda+U_jD_{r+1})\notag\\
 &\quad\left.+\overline U^r\left(
 a_j^2F_r/\lambda^2+2|a_j|U_jE_{r+1}/\lambda+
 U_j^2D_{r+2}+U_{jj}D_{r+1}\right)\right].
 \label{eq:iv-R-second}
\end{align}
The first term in \eqref{eq:iv-R-second} is zero when $r=1$; a negative
power at $\overline U=0$ is not evaluated.  To derive these formulae,
differentiate $U^rQ$ once and twice.  The factors $\partial_jt=z\partial_jU$
and $\partial_{jj}t=z\partial_{jj}U$ are bounded by $U_j$ and $U_{jj}$,
while $\partial_jm=a_j$.  The normalized Beta weight integrates to one.
Every term in the twice-differentiated product appears explicitly above.
The bounds $U_j=\sum_iw_i|s_{ij}|e^{d_i+L\|s_i\|_1}$ and
$U_{jj}=\sum_iw_is_{ij}^2e^{d_i+L\|s_i\|_1}$ are valid on $[-L,L]^p$.

For $f(y)=R_r(y)\exp(-\|y\|^2/2)$,
\begin{equation}
 |\partial_{jj}f|\leq
 \|\partial_{jj}R_r\|_\infty+
 2L\|\partial_jR_r\|_\infty+(L^2+1)\|R_r\|_\infty.
 \label{eq:iv-density-second}
\end{equation}
This follows from the product rule and
$\partial_j\omega=-y_j\omega$,
$\partial_{jj}\omega=(y_j^2-1)\omega$, $0\leq\omega\leq1$.
For an interval of length $2L$, midpoint integration with $N$ equal cells
has expectation error at most $(2L)^2\|f''\|_\infty/(24N^2)$.
One proof integrates the second-order Taylor formula around each cell midpoint;
the linear term integrates to zero and the quadratic remainder integrates
to at most $\|f''\|_\infty h^3/24$.  Summing and dividing by $2L$ gives
the stated expectation bound.  A telescoping tensor-product argument,
applying one coordinate quadrature at a time and using the fact that both
integration and positive midpoint averaging are sup-norm contractions, gives
\begin{equation}
 \beta_{Y,\mathrm{box}}\leq
 C_L\frac{(2L)^2}{24N^2}\sum_{j=1}^p
 \|\partial_{jj}f\|_\infty.
 \label{eq:iv-box-midpoint}
\end{equation}
This grid-only bound still requires a continuous tail certificate.  Its
conservative sufficient widths may be impractical and are retained as failures
rather than assumed to be attainable.

For ideal conditional-normal quantile midpoints, a physical-coordinate
Lipschitz constant $L_j$ gives a grid-only bound $L\sum_jL_j/N$.
To verify the one-coordinate coupling, split each probability cell at its
midpoint.  Monotonicity of the quantile function bounds the two half-cell
integrals by their half-cell probability width times the two corresponding
endpoint increments.  Summing telescopes to
$(1/(2N))(L-(-L))=L/N$.  The multivariate bound follows by successive
coordinate coupling.  Actual dyadic node rounding adds
$\sum_jL_j\max_k|\widehat y_k-y_k|$ separately.  A grid-only width derived
from $L\sum_jL_j/N$ is consequently not a complete financial certificate.

\subsubsection{Complete monetary allowance and precision interpretation}

Let $\beta_Y$ be a certified continuous-to-finite financial transfer,
$e_c$ the radius of the retained coefficient interval, $[F_-,F_+]$ the
external-factor interval, $\overline F=(F_-+F_+)/2$ and
$e_F=(F_+-F_-)/2$.  The implementation uses the interval midpoints.
Let $B=B_2$ for either second correction and $B=B_1$ for the first.
Let $e_{\mathrm{arith}}=2B_2\eta_p$ or $B_1\eta_p$ as appropriate, and
put $\beta_{\mathrm{aux}}=\beta_Z+\beta_S$ for the integral method and zero
for either direct method.  The deterministic undiscounted price error is at most
\begin{equation}
 e_{\mathrm{det}}\leq
 e_c+\beta_Y+F_+(\beta_{\mathrm{aux}}+e_{\mathrm{arith}})+e_FB.
 \label{eq:iv-total-budget}
\end{equation}
To prove this, insert the exact continuous coefficient, exact finite correction,
ideal finite auxiliary average and implemented average one at a time, and
apply the triangle inequality.  The factor approximation can be charged
against the implemented bounded score, whose absolute value is at most $B$;
this avoids an unjustified bound on an unclamped approximate residual.
If a statistical estimator contributes $e_{\mathrm{stat}}$ before the factor,
add $F_+e_{\mathrm{stat}}$.  The final financial price is
$\Pi=e^{-r_f\tau}V$, so the same positive discount factor multiplies the
entire undiscounted bound.  The deterministic and statistical allocations must
sum to the requested monetary tolerance at the stated joint confidence.
A failed allocation is a failed configuration, even if every finite circuit
identity is exact.

The integral factorisation also identifies a local precision consequence.
On a fixed retained support, assume a centred weak-dispersion family with
$U=O(\delta^2)$ and $U_*=\Theta(\delta^2)$, with $H,\lambda$ and support
fixed.  Then $B_2=\Theta(\delta^4)$.  An $h$-table error $\eta_h$ enters
the unnormalised price with coefficient at most
$HU^2D_2/2=O(\delta^4)$; an absolute $U$ error $\eta_U$ enters the smooth
factorised correction through a term of order $O(\delta^2\eta_U)$ plus
higher-order terms on its guarded domain.  These orders follow from
\eqref{eq:iv-product-identity}, the bounded kernel derivatives and
$\partial_U(U^2Q)=2UQ+U^2\partial_UQ$.
A construction which evaluates the three raw capped-price terms separately
has individual absolute table errors which are not suppressed by $\delta^4$
before their final subtraction.  This explains a possible difference between
the precision requirements of the specific retained constructions.  It is
not a theorem against an optimally stable direct implementation: a shared
Taylor kernel, divided difference or suitable conditional integral can retain
the same cancellation.  Neither this local argument nor the finite arithmetic
certificate establishes a novel algorithm or physical quantum advantage.
